\par\medskip\noindent
1.3.8. Sea $p:P\longrightarrow S$ un morfismo de esquemas que sea
\underline{coherente} (es decir, cuasicompacto y cuasiseparado) y
\underline{plano}, y supongamos que
\begin{equation}\tag{1.3.8.1}
 p_*(\mathcal O_P)\xleftarrow{\ \sim\ }\mathcal O_S.
\end{equation}
Como la formación de $p_*(\mathcal O_P)$ conmuta con todo cambio de
base plano (por la coherencia de $p$), las condiciones enunciadas
sobre el esquema relativo $P/S$ son estables por productos cartesianos
y, por tanto, se satisfacen en particular para las potencias
cartesianas $P\times P$, $P\times P\times P$, etc. También son
estables por todo cambio de base plano.

Sea ahora $G$ un esquema \underline{afín} sobre $S$. La relación
(1.3.7.2) se deduce inmediatamente de (1.3.8.1), pues
\[
 \operatorname{Hom}_S(P,G)
 =\operatorname{Hom}_{\mathrm{Alg}}(\underline A,p_*(\mathcal O_P))
 \simeq\operatorname{Hom}_{\mathrm{Alg}}(\underline A,\mathcal O_S),
\]
donde $\underline A$ es el álgebra cuasicoherente sobre $S$ tal que
$G=\operatorname{Spec}(\sisi{A}{\underline A})$. Por la observación
precedente, la misma conclusión se aplica a $P\times P$,
$P\times P\times P$, etc. Asimismo, si $G$ es plano sobre $S$, la
misma conclusión se aplica al esquema relativo $P_G$ sobre $G$
respecto del esquema relativo $G_G$ sobre $G$.

Cuando $P$ y $G$ son esquemas en grupos sobre $S$, y se trabaja con
una topología sobre $(\mathrm{Sch})_{/S}$ menos fina que la topología
canónica, de modo que $P$ y $G$ puedan identificarse con haces sobre
este sitio, todo morfismo de $P$, $P\times P$,
$P\times P\times P$, etc., a $G$ compatible con las secciones
marcadas es trivial. De manera análoga, si $G$ es plano, todo
$S$-morfismo $P\times G\longrightarrow G$ compatible con las secciones
marcadas es trivial, como se ve interpretándolo como un $G$-morfismo
de $P_G$ a $G_G$. Así, se satisfacen las hipótesis de 1.3.5 y también
las de 1.3.6 cuando $G$ es plano. Además, las tres categorías que
intervienen en (1.3.4.1) son rígidas.

%%%%%%%%%%  scan idx 165  |  folio 154  %%%%%%%%%%

1.3.8.2. El caso más interesante en que se satisface (1.3.8.1), y
sigue satisfaciéndose después de toda extensión de la base, es aquel
en que $P$ es un \underline{esquema abeliano} sobre $S$; entonces se
satisfacen incluso las condiciones de 1.3.7. En este caso, las
conclusiones de \sisi{3.5}{1.3.5} y de \sisi{3.6}{1.3.6} son, por
tanto, válidas para todo grupo afín $G$ sobre $S$. Este es el caso que
J.~P. Serre consideró originalmente, y los desarrollos precedentes se
inspiraron manifiestamente en él. Un caso interesante en que se
satisface (1.3.8.1) sin que se satisfaga universalmente es aquel en que
$P$ es el \underline{modelo de Néron} de un esquema abeliano sobre el
cuerpo de fracciones de un anillo de valoración discreta de espectro
$S$.
